\cleardoublepage

\section{绪论}

\subsection{研究背景与意义}
视觉语言导航（Vision-Language Navigation, VLN）是具身智能（Embodied AI）领域的核心任务之一，旨在要求智能体在一个未知的三维环境中，根据自然语言指令进行导航并到达指定目标位置。随着智能机器人技术在服务、物流及安防等领域的深入应用，如何在非结构化、动态且复杂的真实室内环境中实现高精度的指令遵循与路径规划，已成为学术界与工业界关注的热点。

近年来，随着多模态大语言模型（Vision-Language Models, VLM）的蓬勃发展，智能体的跨模态感知与语义对齐能力得到了显著提升。然而，在真实场景的实际应用中，传统的视觉语言导航仍面临若干严峻挑战：
\begin{enumerate}
    \item \textbf{仿真与现实的视场角差异（Field-of-View Gap）}：当前主流的VLN基准（如R2R、RxR）多基于Habitat或Matterport3D仿真器，其提供的12目全景观测（12-View Panoramic Observation）在真实部署时难以复制。真实机器人（如本课题使用的Seeker四目相机）受限于硬件成本与计算功耗，往往只能提供有限的视角视野。这种“观测偏移”会导致模型在仿真中训练出的特征提取能力在实物平台上失效。
    \item \textbf{导航决策的逻辑缺失与不确定性}：大多数端到端导航模型采用直接预测动作的方式，缺乏显式的推理过程。在面对长程路径或复杂指令时，模型极易出现“幻觉”或模式坍塌，无法在复杂环境中进行空间逻辑推理。
    \item \textbf{强化学习训练的不稳定性}：传统的强化学习算法（如PPO）在处理具有长时程依赖和稀疏奖励的导航任务时，往往面临高方差、显存开销大以及收敛速度慢的问题，限制了模型在大规模具身数据集上的训练效率。
\end{enumerate}

针对上述问题，本课题旨在构建一个端到端的视觉语言导航系统，命名为 \textbf{VLN-R1}。该系统受DeepSeek-R1等大模型推理技术的启发，创新性地引入了思维链（Chain-of-Thought, CoT）推理机制与组相对策略优化（Group Relative Policy Optimization, GRPO）算法，旨在提升机器人在复杂环境下的逻辑推理能力与训练稳定性。同时，本课题针对真实硬件Seerker相机设计了4-View特征适配方案，打通了从仿真训练到实地部署的技术通路。

\subsection{主要研究内容}
本文的主要研究工作围绕视觉语言导航系统的感知增强、决策推理及训练优化展开，具体内容如下：
\begin{enumerate}
    \item \textbf{面向实车部署的4-View视觉特征适配}：针对仿真环境12目观测与实车4目观测的差异，提出了一种基于DINOv2特征提取与坐标系对齐的感知迁移方案。通过集成FoundationStereo深度估计模型，实现了在有限视场角下的三维环境鲁棒感知。
    \item \textbf{具有推理能力的CoT具身导航数据集构建}：利用大规模视觉语言模型生成包含空间推理过程的导航思维链（CoT），构建了VLN-Ego数据集。该数据集通过“思维-行动”（Think-Before-Action, TBA）格式，显式地赋予了智能体在决策前的逻辑推理能力。
    \item \textbf{基于GRPO的长程导航策略优化}：将组相对策略优化算法引入具身导航领域，通过组内样本的相对优势估计代替复杂的价值网络训练，显著提升了策略梯度下降的稳定性，解决了长序列动作生成中的训练难点。
    \item \textbf{虚实结合的系统验证与部署}：在R2R-CE与RxR-CE等标准连续环境基准上进行实验，并与Nav-R1、OctoNav等主流算法进行对比。最后，在搭载Seeker相机的移动机器人平台上完成了VLN-R1系统的实物部署，验证了算法的实用性。
\end{enumerate}

\subsection{本文结构安排}
本文共分为八章，各章节安排如下：
第一章绪论，阐述论文的研究背景、意义及核心内容。
第二章相关工作，综述多模态学习、视觉语言导航及强化学习领域的最新研究现状。
第三章方法论与设计思路，详细介绍VLN-R1的系统架构，包括预训练、DAgger引导与GRPO微调。
第四章训练与实验结果，展示在仿真环境下的训练表现及与主流模型的对比分析。
第五章实物部分，介绍Seeker相机的硬件配置、视场角对齐及实车部署流程。
第六章小结，对全文工作进行总结并展望未来方向。
第七章参考文献，列出本文引用的国内外核心文献。
第八章致谢，对提供学业与研究支持的专家学者表示感谢。
